\documentclass[12pt]{article}
\usepackage[utf8]{inputenc}
\usepackage[T1]{fontenc}
\usepackage[margin=1in]{geometry}
\usepackage{setspace}
\usepackage{xcolor}
\usepackage[colorlinks=true, urlcolor=blue!60!black, linkcolor=blue!60!black]{hyperref}
\setlength{\parindent}{2em}
\setlength{\parskip}{1.25em}
\renewcommand{\baselinestretch}{1.0}
\begin{document}
\begin{center}
{\LARGE \textit{United States v. Darby}, 312 U.S. 100 (1941)\par}
\vspace{0.8em}
{\large Case Brief\par}
\end{center}
\vspace{1em}

\section{Memory Jogger}

Overruling \textit{Hammer v. Dagenhart} and dismissing the Tenth Amendment as ``but a truism,'' a unanimous Supreme Court upheld the Fair Labor Standards Act of 1938, holding that Congress may prohibit the interstate shipment of goods produced under substandard labor conditions and directly regulate intrastate production that substantially affects interstate commerce (\href{https://www.courtlistener.com/opinion/103442/united-states-v-darby/}{CourtListener}; \href{https://supreme.justia.com/cases/federal/us/312/100/}{Justia}).

\section{Detailed Case Facts}

In 1938, Congress enacted the Fair Labor Standards Act (FLSA), 52 Stat. 1060, 29 U.S.C. \S~201 \textit{et seq.}, to establish nationwide minimum wages and maximum weekly hours for employees ``engaged in [interstate] commerce or in the production of goods for [interstate] commerce'' (\href{https://supreme.justia.com/cases/federal/us/312/100/}{312 U.S. at 108--110}). During the Act's first year of operation, \S~6 mandated a minimum wage of 25 cents per hour, and \S~7 set a maximum workweek of 44 hours without overtime pay of at least one and one-half times the regular hourly rate (\href{https://supreme.justia.com/cases/federal/us/312/100/}{312 U.S. at 110, 125}). Section 15(a)(1) barred the interstate shipment of goods produced in violation of \S\S~6 or 7; \S~15(a)(2) directly prohibited employing workers in the production of goods for interstate commerce in violation of \S\S~6 or 7; and \S~15(a)(5), incorporating \S~11(c), made it unlawful for covered employers to fail to keep wage and hour records prescribed by the Wage and Hour Administrator (\href{https://supreme.justia.com/cases/federal/us/312/100/}{312 U.S. at 110--111}). Under \S~16(a), willful violations carried criminal penalties of up to a \$10,000 fine for a first offense and up to six months' imprisonment, a fine, or both for subsequent convictions (\href{https://supreme.justia.com/cases/federal/us/312/100/}{312 U.S. at 110}).

Appellee Fred W. Darby (operating the F. W. Darby Lumber Company in Statesboro, Georgia) acquired raw timber in Georgia and manufactured it into finished lumber with the intent to ship a large proportion of that lumber to out-of-state customers, and in fact shipped substantial quantities across state lines (\href{https://supreme.justia.com/cases/federal/us/312/100/}{312 U.S. at 111}; \href{https://law.justia.com/cases/federal/district-courts/FSupp/32/734/2245495/}{32 F. Supp. 734, 735--737 (S.D. Ga. 1940)}). Federal prosecutors obtained a multi-count indictment charging Darby with violating FLSA \S\S~15(a)(1), 15(a)(2), and 15(a)(5) by paying lumber mill workers less than 25 cents per hour, employing them beyond the 44-hour weekly maximum without overtime compensation, shipping the resulting lumber in interstate commerce, and failing to maintain required daily and weekly hour records (\href{https://supreme.justia.com/cases/federal/us/312/100/}{312 U.S. at 108, 111}).

\section{Procedural History}

Darby demurred to the indictment in the United States District Court for the Southern District of Georgia, contending that the FLSA exceeded Congress's Commerce Clause authority and violated the Fifth and Tenth Amendments (\href{https://supreme.justia.com/cases/federal/us/312/100/}{312 U.S. at 108, 111}). On April 29, 1940, District Judge William H. Barrett sustained the demurrer and quashed the indictment in its entirety (\href{https://law.justia.com/cases/federal/district-courts/FSupp/32/734/2245495/}{\textit{United States v. F. W. Darby Lumber Co.}, 32 F. Supp. 734, 738 (S.D. Ga. 1940)}). Relying on \textit{Hammer v. Dagenhart}, 247 U.S. 251 (1918), Judge Barrett ruled that manufacturing is not interstate commerce, that an employer's mere intent at the time of manufacture to ship finished lumber out of state did not bring local production within federal control, and that extending federal labor regulation to local lumber milling would destroy the constitutional balance between federal and state power (\href{https://law.justia.com/cases/federal/district-courts/FSupp/32/734/2245495/}{32 F. Supp. at 737--738}).

Because the district court's judgment quashing the indictment rested on the invalidity and construction of a federal statute, the United States took a direct appeal to the Supreme Court under the Criminal Appeals Act, 18 U.S.C. \S~682, and Judicial Code \S~238, 28 U.S.C. \S~345 (\href{https://supreme.justia.com/cases/federal/us/312/100/}{312 U.S. at 108--109}). The Supreme Court heard oral argument on December 19--20, 1940, with Solicitor General Francis B. Biddle arguing for the United States and Archibald B. Lovett for Darby, and issued its decision on February 3, 1941 (\href{https://www.oyez.org/cases/1940-1955/312us100}{Oyez}).

\section{Judicial Votes}

The Supreme Court reversed unanimously, \textbf{8--0}, with no concurring or dissenting opinions (\href{https://www.oyez.org/cases/1940-1955/312us100}{Oyez}; \href{https://supreme.justia.com/cases/federal/us/312/100/}{Justia}).

\begin{itemize}
\item \textbf{Majority (8):} Justice Harlan Fiske Stone (author), joined by Chief Justice Charles Evans Hughes and Associate Justices Owen J. Roberts, Hugo L. Black, Stanley F. Reed, Felix Frankfurter, William O. Douglas, and Frank Murphy (\href{https://www.oyez.org/cases/1940-1955/312us100}{Oyez}).
\item \textbf{Not Participating / Vacancy (1):} Justice James Clark McReynolds retired effective January 31, 1941---three days before the decision was handed down on February 3, 1941---leaving one vacant seat on the nine-member Court (\href{https://www.cambridge.org/core/journals/american-political-science-review/article/constitutional-law-in-19401941-the-constitutional-decisions-of-the-supreme-court-of-the-united-states-in-the-october-term-1940/FD13DE1C1017FD7C2C676714F9E05EC0}{American Political Science Review}; \href{https://en.wikipedia.org/wiki/List_of_United_States_Supreme_Court_cases,_volume_289}{Hughes Court tenure records}). (Although some secondary databases list all nine Justices from the December 1940 argument, Oyez records the eight sitting Justices who decided the case.)
\item \textbf{Concurrences (0):} None.
\item \textbf{Dissents (0):} None.
\end{itemize}

\section{Holding}

\textbf{Judgment.} Reversed. The district court's order quashing the indictment was reversed, and the FLSA's shipment prohibition (\S~15(a)(1)), direct wage-and-hour regulation of production for interstate commerce (\S~15(a)(2)), and recordkeeping mandate (\S~15(a)(5) and \S~11(c)) were upheld as constitutional (\href{https://supreme.justia.com/cases/federal/us/312/100/}{312 U.S. at 111, 126}).

\textbf{Rule Announced.} Under the Commerce Clause (Art. I, \S~8, cl. 3) and the Necessary and Proper Clause (Art. I, \S~8, cl. 18), Congress possesses plenary authority to prohibit the interstate shipment of any goods produced under substandard labor conditions, regardless of Congress's legislative motive or the indirect effect on intrastate manufacturing; \textit{Hammer v. Dagenhart}, 247 U.S. 251 (1918), is expressly overruled (\href{https://supreme.justia.com/cases/federal/us/312/100/}{312 U.S. at 115--117}). Independently, Congress may directly regulate intrastate production of goods intended or expected in the normal course of business to move in interstate commerce, both as an appropriate means to enforce its interstate shipment ban and because substandard labor conditions in such production substantially affect interstate commerce (\href{https://supreme.justia.com/cases/federal/us/312/100/}{312 U.S. at 117--123}). The Tenth Amendment imposes no independent substantive barrier on Congress's use of appropriate means to carry out a delegated power, and the FLSA's wage, hour, and recordkeeping rules satisfy Fifth Amendment due process and notice requirements (\href{https://supreme.justia.com/cases/federal/us/312/100/}{312 U.S. at 123--126}).

\textbf{Holding vs. Dicta.} The core holdings are: (1) Congress may bar interstate shipment of goods made under substandard wages or hours, overruling \textit{Hammer v. Dagenhart}; (2) Congress may directly regulate wages and hours in intrastate production of goods intended for interstate commerce; (3) the Tenth Amendment does not bar otherwise valid exercises of the commerce power; and (4) the FLSA's recordkeeping and minimum-wage/maximum-hour rules comply with the Fifth Amendment (\href{https://supreme.justia.com/cases/federal/us/312/100/}{312 U.S. at 115--126}). By contrast, Justice Stone's sweeping characterization of the Tenth Amendment as ``but a truism'' and his broad assertion that congressional motive in regulating commerce is entirely beyond judicial review function partly as expansive doctrinal statements (or dicta) that later Courts cabined---for instance, by reviving Tenth Amendment anti-commandeering limits in \textit{New York v. United States}, 505 U.S. 144 (1992), and \textit{Printz v. United States}, 521 U.S. 898 (1997), and limiting the substantial-effects prong to economic activities in \textit{United States v. Lopez}, 514 U.S. 549 (1995) (\href{https://www.cardozolawreview.com/popular-sovereignty-tenth-amendment-reese/}{Cardozo Law Review}; \href{https://www.everycrsreport.com/reports/RL32844.html}{CRS Report RL32844}).

\section{Analysis of Opinions}

\subsection{Majority Opinion (Justice Stone)}

Justice Stone structured the opinion around four constitutional questions: the interstate shipment ban, the direct regulation of local production, the Tenth Amendment, and the recordkeeping and Fifth Amendment objections (\href{https://supreme.justia.com/cases/federal/us/312/100/}{312 U.S. at 112--126}).

First, addressing \S~15(a)(1)'s ban on shipping noncompliant lumber across state lines, Stone acknowledged that manufacturing is not itself interstate commerce, but stressed that shipping manufactured goods across state borders ``is indubitably a regulation of the commerce'' (\href{https://supreme.justia.com/cases/federal/us/312/100/}{312 U.S. at 113}). Quoting Chief Justice Marshall in \textit{Gibbons v. Ogden}, 22 U.S. (9 Wheat.) 1, 196 (1824), Stone reaffirmed that the commerce power is complete in itself, acknowledges no limitations other than those prescribed in the Constitution, and encompasses the power to prohibit interstate movement of goods---not merely noxious articles like lottery tickets or diseased cattle, but any articles Congress deems injurious to public welfare (\href{https://supreme.justia.com/cases/federal/us/312/100/}{312 U.S. at 113--115}). Crucially, Stone rejected judicial inquiry into legislative motive: Congress's purpose to prevent interstate commerce from becoming an instrument of unfair competition that spreads substandard labor conditions was a valid policy judgment beyond judicial control (\href{https://supreme.justia.com/cases/federal/us/312/100/}{312 U.S. at 115}). Finding \textit{Hammer v. Dagenhart} irreconcilable with these principles and praising Justice Holmes's ``powerful and now classic dissent,'' the Court expressly overruled \textit{Hammer} (\href{https://supreme.justia.com/cases/federal/us/312/100/}{312 U.S. at 115--117}).

Second, turning to \S~15(a)(2)'s direct wage-and-hour mandate for intrastate production, Stone first construed ``production for interstate commerce'' to include goods manufactured with the intent or expectation, in the normal course of business, that all or part will be shipped out of state, noting that factories rarely segregate interstate and intrastate inventory on the shop floor (\href{https://supreme.justia.com/cases/federal/us/312/100/}{312 U.S. at 117--118}). Invoking \textit{McCulloch v. Maryland}, 17 U.S. (4 Wheat.) 316, 421 (1819), Stone upheld \S~15(a)(2) on two independent grounds: (a) as a Necessary and Proper means to make the \S~15(a)(1) shipment ban effective by stopping substandard production at its source, and (b) independently of the shipment ban, because intrastate production under substandard wages substantially affects interstate commerce by enabling unfair interstate competition---even by small producers whose aggregate competitive effect is substantial (\href{https://supreme.justia.com/cases/federal/us/312/100/}{312 U.S. at 118--123}). Stone expressly declared \textit{Carter v. Carter Coal Co.}, 298 U.S. 238 (1936), limited insofar as it conflicted with this analysis (\href{https://supreme.justia.com/cases/federal/us/312/100/}{312 U.S. at 123}).

Third, Stone dismantled the district court's federalism objection by declaring that the Tenth Amendment ``states but a truism that all is retained which has not been surrendered'' (\href{https://supreme.justia.com/cases/federal/us/312/100/}{312 U.S. at 124}; \href{https://www.law.cornell.edu/constitution-conan/amendment-10/the-tenth-amendment-and-darby}{Constitution Annotated}). Drawing on ratification debates and Joseph Story's \textit{Commentaries}, Stone concluded that the Amendment was merely declaratory and did not curtail Congress's authority to use appropriate means to execute its granted powers (\href{https://supreme.justia.com/cases/federal/us/312/100/}{312 U.S. at 124}).

Fourth, Stone briefly upheld the \S~11(c) recordkeeping mandate as an appropriate incidental enforcement tool and rejected Darby's Fifth Amendment challenges, citing \textit{West Coast Hotel Co. v. Parrish}, 300 U.S. 379 (1937), to confirm that minimum-wage and maximum-hour laws for both men and women satisfy due process, and finding the statute sufficiently definite to provide criminal notice (\href{https://supreme.justia.com/cases/federal/us/312/100/}{312 U.S. at 124--126}).

\subsection{Concurrences and Dissents}

There were no concurring or dissenting opinions. All eight participating Justices joined Justice Stone's opinion in full (\href{https://www.oyez.org/cases/1940-1955/312us100}{Oyez}).

\section{Examples: Future Applications}

\begin{enumerate}
\item \textbf{Same-Side Application (Valid): Federal Workplace Safety Standards for Regional Manufacturers.} Congress enacts a statute barring interstate shipment of manufactured electronics assembled in facilities that violate federal chemical-exposure limits and directly penalizes employers who expose assembly workers to prohibited toxins while producing electronics destined for out-of-state buyers. Under \textit{Darby}, both the shipment ban and the direct regulation of manufacturing conditions are valid exercises of the Commerce and Necessary and Proper Clauses because Congress may close the channels of interstate commerce to goods made under substandard conditions and regulate intrastate production that substantially affects interstate competition (\href{https://supreme.justia.com/cases/federal/us/312/100/}{312 U.S. at 115--123}).

\item \textbf{Same-Side Application (Valid): Mandatory Compliance Recordkeeping for Small Agricultural Processors.} Congress requires small intrastate food-processing plants that sell a portion of their output to interstate distributors to maintain detailed sanitation and wage logs for all plant employees, even though some batches are sold solely within the home state. Under \textit{Darby}, Congress may regulate commingled intrastate production and impose incidental recordkeeping duties on small producers whose aggregate output affects interstate markets (\href{https://supreme.justia.com/cases/federal/us/312/100/}{312 U.S. at 118, 121--125}).

\item \textbf{Opposite-Side Application (Invalid): Federal Ban on Purely Non-Economic Local Conduct Near Schools.} Congress makes it a federal crime merely to possess a firearm within 1,000 feet of a local school without any requirement that the possession be tied to interstate commerce or an economic enterprise. Unlike the commercial lumber manufacturing in \textit{Darby}, mere non-economic gun possession lacks a substantial economic effect on interstate commerce; as the Court later held in \textit{United States v. Lopez}, 514 U.S. 549 (1995), \textit{Darby}'s substantial-effects doctrine does not authorize a general federal police power over non-economic local crime (\href{https://www.everycrsreport.com/reports/RL32844.html}{CRS Report RL32844}).

\item \textbf{Opposite-Side Application (Invalid): Federal Mandate Compelling Uninsured Individuals to Purchase a Product.} Congress compels individuals who are not participating in the interstate healthcare market to purchase health insurance under the Commerce Clause alone. While \textit{Darby} permits Congress to regulate existing commercial production and interstate shipment, the Supreme Court held in \textit{NFIB v. Sebelius}, 567 U.S. 519 (2012), that the Commerce Clause regulates existing commercial activity rather than compelling individuals to enter commerce in the first place (\href{https://www.everycrsreport.com/reports/RL32844.html}{CRS Report RL32844}).

\item \textbf{Fence-Sitter Application (Close Case): Federal Wage Floor for Strictly Intrastate Artisans Buying Out-of-State Raw Materials.} Suppose a three-person furniture workshop in Georgia sells 100\% of its finished tables exclusively to local walk-in customers and never ships across state lines, but buys its varnish and saw blades from out-of-state suppliers. Under \textit{Darby} alone, the workshop is not ``producing goods for interstate commerce'' because it never intends or expects its lumber to cross state lines (\href{https://supreme.justia.com/cases/federal/us/312/100/}{312 U.S. at 117--118}). Yet under \textit{Darby}'s broader aggregation rationale (later expanded in \textit{Wickard v. Filburn}, 317 U.S. 111 (1942), and \textit{Katzenbach v. McClung}, 379 U.S. 294 (1964)), Congress could argue that the aggregate demand for out-of-state supplies and local competition with interstate furniture makers substantially affects interstate commerce, placing the case squarely on the doctrinal boundary (\href{https://www.everycrsreport.com/reports/RL32844.html}{CRS Report RL32844}).
\end{enumerate}

\section{Critique}

\textbf{Originalist and Textualist Critiques.} Originalist and classical-liberal scholars sharply criticize \textit{Darby} for eroding the Constitution's structure of enumerated powers. Scholars such as Randy Barnett and Richard Epstein argue that at the Founding, ``commerce'' meant trade, barter, and exchange of goods across state lines---distinct from antecedent local manufacturing, agriculture, and labor relations (\href{https://www.cambridge.org/core/product/D258358C733DF44CEB17E8F616237005/core-reader}{The Constitution of 1789, Ch. 16}; \href{https://harvardlawreview.org/forum/vol-128/an-unapologetic-defense-of-the-classical-liberal-constitution/}{Epstein, Harvard Law Review Forum}). From this perspective, by abandoning judicial review of pretextual legislative motives (contrary to Chief Justice Marshall's warning in \textit{McCulloch v. Maryland} against Congress using enumerated powers as a ``pretext'' to accomplish unenumerated objects) and collapsing the Tenth Amendment into a mere ``truism,'' \textit{Darby} effectively handed Congress a general national police power (\href{https://conlaw.us/case/united-states-v-darby-1941/}{ConLaw.us}; \href{https://www.heritage.org/the-constitution/report/commerce-commerce-everywhere-the-uses-and-abuses-the-commerce-clause}{Heritage Foundation}). Critics also contend that treating the Tenth Amendment as a nullity ignored its structural role in preserving both state sovereignty and popular sovereignty (``or to the people'') (\href{https://www.cardozolawreview.com/popular-sovereignty-tenth-amendment-reese/}{Cardozo Law Review}).

\textbf{Progressive and Doctrinal Defenses.} Progressive and nationalist scholars defend \textit{Darby} as a necessary correction to the formalistic and economically untenable distinctions of the \textit{Lochner} and \textit{Hammer} eras. Scholars such as Bruce Ackerman view \textit{Darby} as the culmination of the New Deal constitutional transformation, bringing judicial doctrine into alignment with an integrated national economy in which individual states faced a collective-action ``race to the bottom'' if they attempted to enact humane wage, hour, or child-labor standards unilaterally (\href{https://www.law.cornell.edu/constitution-conan/article-6/section-None/clause-2/the-new-deal-and-the-presumption-against-preemption}{Constitution Annotated}; \href{https://supreme.justia.com/cases/federal/us/312/100/}{Justia}). Moreover, revisionist originalist scholarship (such as Andrew Coan and David S. Schwartz) argues that the Founding-era meaning of enumerationism was far more contested and indeterminate than strict enumerationists claim, lending historical plausibility to \textit{Darby}'s deference to Congress on national economic problems (\href{https://ilr.law.uiowa.edu/sites/ilr.law.uiowa.edu/files/2024-03/ILR-109-Coan-Schwartz_0.pdf}{Coan \& Schwartz, Iowa Law Review}).

\textbf{Underlying Competing Values.} At bottom, \textit{Darby} mediates a permanent tension between \textit{national democratic capacity}---empowering Congress to solve interstate collective-action problems like labor exploitation and unfair competition---and \textit{constitutional federalism}, which divides sovereignty between federal and state governments to preserve local self-governance and check centralized power (\href{https://www.law.cornell.edu/constitution-conan/amendment-10/the-tenth-amendment-and-darby}{Constitution Annotated}; \href{https://www.everycrsreport.com/reports/RL32844.html}{CRS Report RL32844}).

\section{Key Quotations}

\begin{quote}
``Whatever their motive and purpose, regulations of commerce which do not infringe some constitutional prohibition are within the plenary power conferred on Congress by the Commerce Clause.''
\end{quote}
\noindent --- Justice Harlan Fiske Stone, Majority Opinion, 312 U.S. at 115 (\href{https://supreme.justia.com/cases/federal/us/312/100/}{Justia}; \href{https://www.courtlistener.com/opinion/103442/united-states-v-darby/}{CourtListener}).

\begin{quote}
``The conclusion is inescapable that Hammer v. Dagenhart was a departure from the principles which have prevailed in the interpretation of the Commerce Clause both before and since the decision, and that such vitality, as a precedent, as it then had, has long since been exhausted. It should be, and now is, overruled.''
\end{quote}
\noindent --- Justice Harlan Fiske Stone, Majority Opinion, 312 U.S. at 116--117 (\href{https://supreme.justia.com/cases/federal/us/312/100/}{Justia}; \href{https://www.courtlistener.com/opinion/103442/united-states-v-darby/}{CourtListener}).

\begin{quote}
``The power of Congress over interstate commerce is not confined to the regulation of commerce among the states. It extends to those activities intrastate which so affect interstate commerce or the exercise of the power of Congress over it as to make regulation of them appropriate means to the attainment of a legitimate end, the exercise of the granted power of Congress to regulate interstate commerce.''
\end{quote}
\noindent --- Justice Harlan Fiske Stone, Majority Opinion, 312 U.S. at 118 (\href{https://supreme.justia.com/cases/federal/us/312/100/}{Justia}; \href{https://www.courtlistener.com/opinion/103442/united-states-v-darby/}{CourtListener}).

\begin{quote}
``The amendment states but a truism that all is retained which has not been surrendered. There is nothing in the history of its adoption to suggest that it was more than declaratory of the relationship between the national and state governments as it had been established by the Constitution before the amendment, or that its purpose was other than to allay fears that the new national government might seek to exercise powers not granted, and that the states might not be able to exercise fully their reserved powers.''
\end{quote}
\noindent --- Justice Harlan Fiske Stone, Majority Opinion, 312 U.S. at 124 (\href{https://supreme.justia.com/cases/federal/us/312/100/}{Justia}; \href{https://www.courtlistener.com/opinion/103442/united-states-v-darby/}{CourtListener}).

\end{document}
